\documentclass[11pt]{article}

\usepackage{amsmath,amssymb,amsthm}
\usepackage[margin=1in]{geometry}
\usepackage{hyperref}

\newcommand{\C}{\mathbb{C}}
\newcommand{\Z}{\mathbb{Z}}
\newcommand{\g}{\mathfrak{g}}
\newcommand{\p}{\mathfrak{p}}
\newcommand{\uu}{\mathfrak{u}}
\newcommand{\orr}{\operatorname{or}}
\newcommand{\sig}{\operatorname{sig}}
\DeclareMathOperator{\dimm}{dim}

\title{The role of orientation numbers in the \texttt{atlas} unitarity algorithm}
\author{}
\date{\today}

\begin{document}
\maketitle

\section{The question}

The \texttt{atlas} algorithm for deciding whether an irreducible representation
$\pi=J(\gamma)$ is unitary proceeds in two stages:
\begin{enumerate}
  \item[(1)] compute the signature of the \emph{$c$-form} on $\pi$;
  \item[(2)] \emph{convert} that signature into the signature of the genuine
        invariant Hermitian form on $\pi$.
\end{enumerate}
The representation is unitary iff the resulting Hermitian form is definite
(\texttt{is\_pure} in the code).  The question is: where do orientation numbers
enter?  Do they belong to stage~(1), the $c$-form computation, or only to
stage~(2), the Hermitian conversion?

\medskip
\noindent\textbf{Short answer.}  Orientation numbers belong to stage~(1).  They
appear in the passage from \emph{standard} modules to the \emph{irreducible}
module \emph{inside the $c$-form computation}.  They do \emph{not} appear in the
$c$-to-Hermitian conversion of stage~(2); that conversion is governed by a
different invariant $\mu$ (built from $\dimm(\uu\cap\p)$ and the value of the
lowest $K$-type).  So the naive guess --- ``orientation numbers are the thing
that turns the $c$-form into the Hermitian form'' --- is, perhaps surprisingly,
\emph{not} how the algorithm is organized.

\section{What an orientation number is}

The signature characters live in the ring of \emph{split integers}
$\Z[s]/(s^2-1)$: the unit $s$ records a sign, i.e.\ it interchanges the positive
and negative parts of a signature $(p,q)$.

The orientation number $\orr(\gamma)$ of a parameter $\gamma=(x,\lambda,\nu)$ is a
nonnegative integer attached to $\gamma$.  It is computed by the built-in
\texttt{orientation\_nr} (the C\texttt{++} routine
\texttt{Rep\_context::orientation\_number}), and counts the \emph{non-integral}
positive roots that contribute, namely
\begin{itemize}
  \item each non-integral positive \emph{real} root $\alpha$ that is
        ``oriented'' with respect to $\gamma$ (a parity condition on
        $\langle\gamma-\lambda+\rho_{\mathbb R},\alpha^\vee\rangle$ together with
        the sign of $\langle\nu,\alpha^\vee\rangle$), and
  \item each non-integral complex quadruple $\{\pm\alpha,\pm\theta\alpha\}$ that
        is a complex descent (these are counted in pairs and then halved).
\end{itemize}
The essential point for us is that $\orr(\gamma)$ depends on the continuous
parameter $\nu$: it is a property of where $\gamma$ sits relative to the
non-integral root hyperplanes, not merely of the block.

\section{Stage (1): the $c$-form, and where orientation numbers live}

\subsection*{$c$-form on a standard module: no orientation numbers}

The $c$-form on a \emph{standard} module $I(x)$ is computed purely by
deformation in $\nu$ (\texttt{c\_form\_std} $=$ \texttt{full\_deform}).  As
$\nu\to0$ one reaches the tempered case, where the $c$-form is positive; the
deformation accumulates the controlled sign changes at reducibility points,
encoded by Kazhdan--Lusztig--Vogan polynomials evaluated at $q=s$.  No
orientation number is involved in this step.

\subsection*{$c$-form on an irreducible module: orientation numbers enter here}

To get the $c$-form on the irreducible $J(y)$ one inverts the
Kazhdan--Lusztig--Vogan relationship.  In the Grothendieck group,
\begin{equation}
  J(y)=\sum_x (-1)^{\ell(x)-\ell(y)}\,P_{x,y}(1)\,I(x).
\end{equation}
This identity of \emph{characters} does \emph{not} carry over verbatim to
\emph{forms}.  At the level of $c$-forms one has instead
\begin{equation}\label{eq:cform-irr}
  J(y)_c=\sum_x (-1)^{\orr(x)-\orr(y)}\,(-1)^{\ell(x)-\ell(y)}\,P_{x,y}(s)\,I(x)_c,
\end{equation}
where two things change relative to the character formula:
\begin{itemize}
  \item the polynomial is evaluated at $q=s$ rather than $q=1$ (this is the
        signature refinement of the multiplicity), and
  \item each term acquires the \textbf{orientation-number sign}
        $(-1)^{\orr(x)-\orr(y)}$.
\end{itemize}
The difference $\orr(x)-\orr(y)$ is always even on terms that occur (an odd
difference signals an error), so this sign is well defined; in the split ring it
is the factor $s^{(\orr(y)-\orr(x))/2}$.

In the code this is exactly \texttt{c\_form\_irreducible}:
\begin{verbatim}
set c_form_irreducible (Param p) = KTypePol:
   let ori_nr_p = orientation_nr(p)
   then oriented_KL_sum = ParamPol: null_module(p) +
      for c@q in KL_sum_at_s(p) do (c*orientation_nr_term(ori_nr_p,q),q) od
in map(c_form_std@Param, oriented_KL_sum )
\end{verbatim}
Here \texttt{KL\_sum\_at\_s(p)} supplies $(-1)^{\ell(x)-\ell(y)}P_{x,y}(s)$, and
\begin{verbatim}
orientation_nr_term(orientation_nr(p),q) = s^{[or(p)-or(q)]/2}
\end{verbatim}
supplies the orientation-number sign.  Only after these standard-module
coefficients are fixed up is each $I(x)_c$ replaced by its deformed value
(\texttt{c\_form\_std}).  The same factor appears at the built-in level as
\texttt{exp\_i(orient\_y - orientation\_number(x))} in
\texttt{Rep\_context}'s deformation routines.

\medskip
\noindent
\emph{Why it must be here.}  The $c$-form is canonically normalized on each
\emph{standard} module by the deformation (its sign is pinned down at the
tempered endpoint).  When one re-expresses $J(y)_c$ in that fixed basis of
standard $c$-forms, the relative normalization between $J(y)$ and each $I(x)$ is
precisely $(-1)^{\orr(x)-\orr(y)}$.  Orientation numbers are therefore an
intrinsic part of computing the $c$-form \emph{on the irreducible}, even though
they are absent from the $c$-form on a single standard module.

\section{Stage (2): converting the $c$-form to the Hermitian form}

The conversion is \texttt{convert\_cform\_hermitian}, and it does \emph{not} use
orientation numbers.  It multiplies the $c$-form's value on each $K$-type $p$ by
$s$ raised to $\mu(p)-\mu(p_0)$, where $p_0$ is a chosen normalizing $K$-type and
\begin{equation}
  \mu(p)= \underbrace{(g-l-\rho^\vee_{\mathbb R})\cdot(\lambda-\rho)}_{\text{$\lambda$-term}}
        + \underbrace{\tfrac12\,l\,(\delta-1)\,\tau}_{\text{$\tau$-term}}
        + \underbrace{\dimm(\uu\cap\p)}_{\text{dimension term}} .
\end{equation}
In code:
\begin{verbatim}
set convert_cform_hermitian (KTypePol P, mat delta, KType p0) = KTypePol:
  let a_mu = mu(p0,delta) in
    P.null_K_module +
      for c@p in P do (c*(mu(p,delta)-a_mu).rat_as_int.exp_s, p) od
\end{verbatim}
The governing quantity here is $\mu$ --- dominated by the parity of
$\dimm(\uu\cap\p)$ together with the lowest-$K$-type valuation --- and not the
orientation number.  This is the only place the genuine real form (via $\delta$
and $\dimm(\uu\cap\p)$) enters the sign bookkeeping; in the equal-rank case the
twist is trivial and \texttt{hermitian\_form\_irreducible} is just
\texttt{c\_form\_irreducible} followed by this conversion.

\section{Conclusion}

\begin{center}
\begin{tabular}{lll}
\hline
Step & atlas routine & orientation numbers?\\
\hline
$c$-form on standard $I(x)$           & \texttt{c\_form\_std}\,/\,\texttt{full\_deform} & no \\
$c$-form on irreducible $J(y)$        & \texttt{c\_form\_irreducible}                   & \textbf{yes} \\
convert $c$-form $\to$ Hermitian      & \texttt{convert\_cform\_hermitian}              & no (uses $\mu$)\\
\hline
\end{tabular}
\end{center}

\noindent
Orientation numbers are part of \textbf{computing the $c$-form}, specifically the
step that assembles the $c$-form on the irreducible from the $c$-forms on
standard modules via formula~\eqref{eq:cform-irr}.  They are \textbf{not} the
mechanism that converts the $c$-form into the Hermitian form; that role is played
by $\mu$ and, in particular, by $\dimm(\uu\cap\p)$.

\end{document}
